% Generated draft: requires mathematical and source review.
\documentclass[11pt]{article}
\usepackage[margin=25mm]{geometry}
\usepackage{fontspec}
\setmainfont{Latin Modern Roman}
\usepackage{amsmath,amssymb,mathtools,unicode-math}
\setmathfont{Latin Modern Math}
\usepackage{enumitem,needspace,xurl,hyperref}
\setlength{\parindent}{0pt}
\setlength{\parskip}{5pt}
\setlength{\emergencystretch}{3em}
\newcommand{\N}{\mathbb{N}}
\newcommand{\Z}{\mathbb{Z}}
\newcommand{\Q}{\mathbb{Q}}
\newcommand{\R}{\mathbb{R}}
\newcommand{\C}{\mathbb{C}}
\newcommand{\F}{\mathbb{F}}
\newcommand{\A}{\mathbb{A}}
\newcommand{\PP}{\mathbb{P}}
\newcommand{\E}{\mathbb{E}}
\newcommand{\eps}{\varepsilon}
\newcommand{\dd}{\,\mathrm{d}}
\begin{document}
\section*{UnsolvedMath ID: 30005288}
\textbf{Source problem code:} catalogue entry
\subsection{Definitions and mathematical statement}
A subset of a set \(M\) is cocountable if its complement in \(M\) has cardinality at most \(\aleph_0\). Under the usual definability convention, a structure is quasiminimal if every subset of its universe definable by a formula in one free variable, allowing parameters from the structure, is countable or cocountable. The supplied record does not specify the permitted language, its cardinality, or the precise conventions underlying its terminology.

Write \(\operatorname{QM}(M)\) and \(\operatorname{QME}(M)\) for the source's assertions that \(M\) is quasiminimal and quasiminimally excellent, respectively. The record supplies no axiomatization of quasiminimal excellence, no associated closure operator, and no definition of its additional abbreviation \(\operatorname{QMAut}\). These predicates therefore retain their source meanings; no alternative axiomatization is imposed here.

The question, identified as Question 5 in the supplied excerpt, is whether
\[
\forall M\;\Bigl(\bigl(|M|\geq\aleph_2\ \land\ \operatorname{QM}(M)\bigr)
\Longrightarrow \operatorname{QME}(M)\Bigr),
\]
where the quantifier ranges over structures in the source's intended setting. Thus the cardinality bound includes \(\aleph_2\) itself and every larger cardinal. The missing specifications of that setting and of \(\operatorname{QME}\) prevent this draft from being a self-contained axiomatization of the question.

The supplied curation records the implication \(\operatorname{QME}\Rightarrow\operatorname{QM}\). The excerpt also presents an \(\aleph_1\)-sized counterexample below the proposed bound: the dense linear order
\[
D=\omega_1\times_{\mathrm{lex}}\Q,
\qquad
(\alpha,q)<(\beta,r)
\iff \alpha<\beta\ \lor\ (\alpha=\beta\ \land\ q<r),
\]
which it describes as \(\operatorname{QMAut}\) but not \(\operatorname{QME}\). The excerpt contrasts approximation of this order by countable substructures along \(\omega_1\) with approximation in the excellent case along the directed partial order of all countable subsets of an uncountable basis, ordered by inclusion. It suggests that work of Pillay and Tanovic, cited there as [PT11], may be relevant. These are supplied contextual claims, not additional hypotheses or independently verified results.

As further context, the excerpt attributes the following statement to Gallinaro and Kirby as Theorem 2, marked ``forthcoming'' at the time of that source. For each \(\lambda\in\C\), put
\[
\Gamma_\lambda=\{(\exp(z),\exp(\lambda z)):z\in\C\},
\]
the graph of the multivalued complex power map \(w\mapsto w^\lambda\). The excerpt states that
\[
\langle\C;+,\cdot,-,0,1,(\Gamma_\lambda)_{\lambda\in\C}\rangle
\]
is quasiminimal. This attributed contextual statement is separate from the question above; its historical publication qualification is retained without a current-status assertion.
\subsection{Short English statement}
Must every quasiminimal structure of cardinality at least \(\aleph_2\) be quasiminimally excellent? The supplied record labels this question open and describes a counterexample at cardinality \(\aleph_1\), but leaves the precise excellence axioms and ambient language conventions unspecified.
\subsection{Sources}
\begin{itemize}
\item Local UnsolvedMath registry: \texttt{data/unsolvedmath/problems.json}, with additional registry records in \texttt{catalog/problems.json}; permanent ID \(30005288\), source code \texttt{catalogue entry}, title ``Quasiminimal Excellence Above Aleph-Two.'' This draft uses the supplied complete record; the registry files were not independently inspected.
\item Supplied report citation: \emph{Mini-Workshop: Topological and Differential Expansions of o-minimal Structures} (2023), \emph{Oberwolfach Reports}, pp.\ 3005--3050, DOI \texttt{10.4171/owr/2022/52}. The excerpt identifies Question 5 and carries page number 3037. \url{https://doi.org/10.4171/owr/2022/52}. The registry records proposed year 2022; the supplied citation uses 2023.
\item Kirby, \emph{Around Zilber's quasiminimality conjecture}, supplied literature reference: \url{https://arxiv.org/abs/2306.05811}.
\item EMS Press, OWR 2022/52, supplied literature reference: \url{https://doi.org/10.4171/owr/2022/52}.
\item Literature assessment dated exactly \texttt{2026-08-22}: the supplied curation labels the question open and attributes to Kirby the implication from quasiminimal excellence to quasiminimality and the \(\aleph_1\) counterexample. This is supplied dated curation derived from the OWR record, not an independently verified literature assessment. The supplied registry also labels the statement \texttt{corrected\_verified} and says it was reconstructed and checked against the official report; that verification claim is likewise reported rather than independently established here.
\item The excerpt mentions Pillay and Tanovic [PT11], without a full bibliographic entry or link, and attributes its contextual Theorem 2 to Gallinaro and Kirby, marked ``forthcoming.'' No further bibliographic information for either attribution is supplied.
\end{itemize}
\end{document}
